% Proposed research extension, 2026-11-24. Included after the existing appendices.
\section{Cashu: reserve capital, claim-once redemption, and exit bandwidth}
\label{app:cashu-capital}
\paragraph{Scope.}
The proof-admission and Cashu wallet prototype in \texttt{proof-grove/} supplies a further candidate settlement object. A Cashu note is a claim on a mint under the mint's denomination-specific blind-signature key; a normal recipient can ask the mint to swap it, and the mint maintains the spent-secret state. The cryptographic model and state interface are specified in NUT-00 and NUT-07.\footnote{\url{https://cashubtc.github.io/nuts/00/}; \url{https://cashubtc.github.io/nuts/07/}.} A valid blind signature says nothing by itself about a reserve output or a unilateral exit. The object studied below is a proposed \emph{protected} profile with an explicitly admitted claim and an enforceable recovery right. It is not the standard transferable bearer note. This distinction is material to the authority and exit components of the tuple in Section~\ref{sec:ark}.

\subsection{The reserve bound and its clock}

Fix a chain, a confirmation policy, a set of protected claims, an adversary model, and a settlement horizon. At time $t$, write $L_t=\sum_{i\in A_t}a_i$ for the face value of the outstanding claims in the authenticated active set $A_t$, in satoshis. Let $R_t$ be the sum of distinct, eligible, unspent reserve outputs that the settlement mechanism can apply to these claims under the stated adversary model. An output pledged elsewhere is excluded unless its competing encumbrances are subtracted. Assets that need the adversary's voluntary signature for recovery are excluded when the adversary is allowed to refuse. Counting the same output through two applications or two roots does not create a second reserve. Let $F_t$ be the additional fees the contract promises to pay under the designated joint redemption scenario, excluding fees borne by the claimants themselves. $F_t$ is a budget under a named fee scenario, not an unconditional upper bound on future fee-market prices.

\begin{proposition}[protected reserve capital]
\label{prop:cashu-reserve}
Suppose all outstanding protected claims may simultaneously exercise their face-value payout rights against the same closed reserve, with no external top-up or future receipts, and all promised settlement fees must be paid by that reserve. Then
\begin{equation}
R_t\geq L_t+F_t.
\label{eq:cashu-reserve}
\end{equation}
If recovery pays net of claimant-borne fees, the corresponding necessary bound uses the net contractual payout plus the fees actually deducted from the reserve; a face-value guarantee must not be inferred from that contract.
\end{proposition}

\begin{proof}
Take the admissible execution in which every outstanding claim exercises its payout right. The sum of the required outputs is $L_t$, the fees chargeable to the reserve are $F_t$, and conservation of Bitcoin inputs and outputs bounds their sum by the distinct eligible inputs available at $t$. Change and other nonnegative outputs cannot reduce the required inputs. This is a necessary condition; the existence of sufficient value alone supplies neither valid transaction paths nor timely confirmation.
\end{proof}

Equation~\eqref{eq:cashu-reserve} is a stock constraint. It differs from the directional prefix condition of Lemma~\ref{lem:cut}: rapid off-chain turnover can reuse a funded stock, while no policy creates reserve principal for simultaneously exercisable claims. At fixed payment volume and horizon, with $V_H>0$ the value of successfully delivered application payments (counted once, excluding internal swaps and reserve movements that would double-count service), apply the capital-time accounting of Section~\ref{sec:ratio}:
\begin{equation}
\mathcal D_{\mathrm{Cashu}}(H)
 =\frac{\int_0^H R_t\,dt}{V_H}
 \geq \frac{\int_0^H (L_t+F_t)\,dt}{V_H}.
\label{eq:cashu-duration}
\end{equation}
The inequality holds pointwise on histories satisfying the protected reserve premises. At zero delivered volume the ratio is undefined. Bond capital locked in separate dispute contracts and additional Lightning liquidity are separate terms in the total numerator; a unit of Bitcoin cannot serve two incompatible obligations at once. Value released to the recipient ceases to be reserve capital, whereas change still committed to a successor state remains in the integral. The fee floor and the lock duration therefore matter even when the proof itself is small.

Partition liabilities among $m$ isolated holders with disjoint reserves, $R_j\geq L_j+F_j$. Pooling into a single reserve gives the necessary condition $R_{\mathrm{pool}}\geq\sum_j L_j+F_{\mathrm{pool}}$. The principal floor is identical. If a realizable batched settlement satisfies $F_{\mathrm{pool}}<\sum_jF_j$, the pooled floor is smaller by the fee-budget difference. This comparison concerns the floors; it is not a claim that a particular pooling protocol attains either floor. A pooled recovery game may require more fee headroom or additional operator and challenger bonds. Forecast-based fractional backing changes the guarantee in the proposition and must be evaluated under a different insolvency and run model.

\subsection{A keys-alone baseline and a transfer counterexample}

For a single holder, fund a SegWit v0 P2WSH output at admission with two script branches: (i) a two-of-two signature condition between holder and mint; (ii) the holder's signature after a relative block delay $d$. The second branch uses BIP 68 and BIP 112 with a compatible transaction version and input sequence.\footnote{\url{https://github.com/bitcoin/bips/blob/master/bip-0068.mediawiki}; \url{https://github.com/bitcoin/bips/blob/master/bip-0112.mediawiki}.} This uses deployed script and does not require a new covenant. The holder must obtain and verify the funding outpoint, exact script, key, amount, chain observation and recovery transaction construction before considering issuance protected. Cooperative redemption requires a transaction the holder agrees to sign; neither party alone can spend through that branch.

\begin{proposition}[whole-output exclusion]
\label{prop:cashu-exclusion}
For a confirmed protected reserve outpoint $o$, suppose the entire cooperative redemption and the entire unilateral recovery each consume $o$. On any valid Bitcoin chain at most one of these spends can be confirmed. Before the relative lock matures, the unilateral branch is unavailable.
\end{proposition}
\begin{proof}
Bitcoin input validity permits each outpoint to be spent at most once on a single valid chain. The unilateral script branch additionally requires the relative age condition. A reorganization can replace which spend confirms, so final application status must track the active chain.
\end{proof}

The result concerns competing Bitcoin spends of the same collateral. It does not cover an ordinary Lightning melt paid from the mint's other wallet, a partial melt that leaves the old outpoint intact, duplicate blind issuance against the same output, or a claimant who sells an old note after recovering the reserve. An authenticated registry must reserve the funding outpoint once and protected minting and melting must use the registry's rules. A holder can still retain a byte copy of a spent note. A recipient must verify the protected claim's current authority and unconsumed status, rather than accept the signature alone. A mint-controlled key can sign arbitrarily many ordinary notes outside this profile; the protected class needs independently verifiable admission and recipients must refuse unadmitted notes. The example assumes holder key custody, active-chain monitoring and a viable fee budget; the time lock gives an earliest valid spend, not a latest confirmation time.

A counterexample prevents an easy transplant of Cashu's usual transfer semantics. Alice holds an output with the delayed branch to her key and sends an associated blind-signed note to Bob. Even if the mint's database marks Alice's note spent and issues Bob a new one, Alice can exercise the old branch when it matures while Bob holds the purportedly protected note. Changing a nullifier in an off-chain root cannot revoke a consensus-valid old signature path. A protected transfer must consume and replace that path, give the new owner a genuinely enforceable revocation/contest mechanism, or settle into a different construction whose current state governs the exit. For the first experiment, the protected claim is a holder-bound, indivisible redemption lot, with whole-lot on-chain cooperative redemption or delayed recovery. General bearer transfers, arbitrary partial melts and Lightning melts are outside this baseline. Rekeying by spending the whole output to a new confirmed reserve resets its age. Privacy of ordinary blind issuance does not imply unlinkability of the public reserve or of such a transfer.

\subsection{Recursive state and enforcement}

The authenticated transition state may commit to $(\mathcal R_t,\mathcal L_t,\mathcal N_t,\mathcal I_t)$, the reserve set, outstanding protected liabilities, consumed claim identifiers, and admitted issuance intents. A candidate relation requires authenticated membership and nonmembership, unique reserve assignment, bounded integer sums, domain-separated encodings, authorization, exact state-root continuity, and the invariant in Equation~\eqref{eq:cashu-reserve} under its fee policy. Issuance admits a new funding outpoint once and binds the claim, amount, denomination/keyset, protected ownership conditions, and the exact blinded issuance request; a swap consumes old claims and creates admissible successors without inflating liability. Redemption consumes its claims once, removes the spent reserve inputs, admits eligible change if any, and checks the payout and fees from exact Bitcoin transactions. Multiple claim identifiers may be consumed in a single transaction; each must be distinct. A public nullifier cannot be an arbitrary caller-provided string; it must be bound to the authenticated claim and its authorization. At a minimum, the witness relation must fix the claim identifier before spending and prove its linkage to the issued entitlement. A scheme deriving identifiers from Cashu secrets must demonstrate how it preserves blindness and prevents alternate encodings or reuse across keysets. Publishing a stable per-note identifier or reserve link may reduce privacy; measure this explicitly.

This is analogous to the unique-exposure relation and epoch restriction in Section~\ref{sec:proofs}, but the object here is a monetary liability. One proof of a valid transition can establish consistency of a proposed state. It cannot establish that the mint has accounted for every liability it privately issued, that a different root was not used elsewhere, or that the settlement output is enforceably bound to the accepted root. A verifier must check the pinned initial state, monotonic successor selection, chain context, and the set of protected signing/admission rules; a relying recipient needs the data and authenticated path necessary to evaluate its own claim. For an optimistic BitVM3 settlement, the disputed relation, collateral-spend transaction graph, challenger access to witnesses or commitments, timeout, operator and challenger incentives, fees, and eventual payout must all be specified separately (cf. Section~\ref{sec:bond}). A successful verifier call is not a collateral contract. Until those obligations are discharged, pooled off-chain transferable protected Cashu is a research target.

\subsection{A second bound: simultaneous recovery bandwidth}

Let $n$ holders require separate on-chain payouts to distinct designated P2WPKH destinations after a common refusal, and assume each destination must receive an output confirmed within a window of $W$ blocks. Let $b$ denote the guaranteed virtual-byte budget per block available to this recovery workload, conditional on fee and inclusion assumptions. A P2WPKH output serializes to $8+1+22=31$ non-witness bytes. Ignoring inputs, transaction headers, dispute transactions and all other outputs gives the necessary bound
\begin{equation}
 W\ \geq\ \left\lceil\frac{31n}{b}\right\rceil,\qquad b>0.
\label{eq:cashu-bandwidth}
\end{equation}
This remains a lower bound if the reserve uses a single pooled output or an aggregate proof. It assumes distinct direct Bitcoin outputs; a different rail, shared destination, or deferred exit changes the requirement. A common-output serial update path may have further contention as in Section~\ref{sec:throughput}; an independently spendable transaction tree can remove that particular serial dependency but must still establish data availability, valid paths, and inclusion capacity. No fixed $b$ or upper recovery time is guaranteed by Bitcoin. At $b=0$ the stipulated on-chain deadline is infeasible for $n>0$.

\subsection{Evidence boundary and proposed tests}

The archived November 23--24, 2026 Codex campaign records a paired regtest Shinigami run proving 25,000 sat issuance in 241 seconds and a 24,750 sat offline payout in 258 seconds, at a reported combined cloud charge of USD~0.2487. A separate cooperative melt established a 24,000 sat holder reserve and a 23,750 sat payout with the mint offline; Cairo execution matched 24 public output words, and a fresh proof for that exact melt-funded path was still outstanding at the last applicable report. These are historical reports in \texttt{codex-campaign-2026-11-23/sessions/}, not newly reproduced measurements in this paper. The initial melt in the latter experiment required mint cooperation. The funding proof, payout proof, recursive state proof and funded enforcement are different evidence obligations. The existing paper's simulated-proof results retain their original scope.

Proposed test PG-AS-01 begins with an isolated regtest chain, a funded individual reserve and a protected issuance record. Test correct and premature recovery, mint-only spend refusal, duplicate issuance, competing cooperative and unilateral spends in both orders, a funding reorganization, changed holder key, stale transferred note, unauthorized mint issuance and fee-bump behavior. Record raw transactions, scripts, block heights, fees, public statement and verifier identifiers and independent Core results. Use no mainnet keys or funds.

PG-AS-02 freezes the same issuance, transfer and redemption histories for disjoint versus pooled reserves, with holder counts $1,8,32,128,512$, identical face amounts and a declared fee-rate grid. Report reserve principal, fee budgets, separate bonds, capital-time per unit delivered, actual transaction vbytes, recovery latency, verifier/prover cost and peak memory. Include synchronized refusal, staggered ordinary redemptions, a reorganization and a missing-witness challenge. A reference state-machine test can falsify a claimed invariant but cannot certify a script or deployed proof verifier. The pooled arm is admitted only after a valid transaction and dispute construction exists; otherwise record it as an unimplemented target. The main comparison is whether any saved fee or forecast capital exceeds the added capital-time and exit costs. No experiment in this appendix has been run for this revision.
